%------------------------------------------------------------------------------%
\documentclass[fleqn,utf8,aspectratio=169,12pt]{beamer}

\usepackage{integracao_e_series}

\title{Integrais Definidas}

%------------------------------------------------------------------------------%
\begin{document}

\addtitle

%------------------------------------------------------------------------------%
\section{Integrais Definidas}

%------------------------------------------------------------------------------%
\begin{frame}{Fórmulas para o cálculo de áreas}
  Figuras regulares

  \pause
  \bigskip
  Área de um Retângulo
  \tabto{45mm}
  \(
    A = (\text{base}) \times (\text{altura}) = bh
  \)

  \pause
  \bigskip
  Área de um Triângulo
  \tabto{45mm}
  \(
    A = \frac{(\text{base}) \times (\text{altura})}{2} = \frac{bh}{2}
  \)

  \pause
  \bigskip
  Área de um Círculo
  \tabto{45mm}
  \(
    A = \pi (\text{raio})^2 = \pi r^2
  \)

  \pause
  \bigskip
  O que fazer com as figuras irregulares?
\end{frame}

%------------------------------------------------------------------------------%
\begin{frame}{Área entre o gráfico de uma função e o eixo $x$}
  \centering
  \bigskip
  \input{fig_integral_definida_area}
  \pause
  \emph{Como calcular a área?}
\end{frame}

%------------------------------------------------------------------------------%
%------------------------------------------------------------------------------%
\begin{question}
  Calcule a área abaixo do gráfico da função \(f(x) = x + 2\) no intervalo
  \([0, b]\)
\end{question}

%------------------------------------------------------------------------------%
\begin{resolution}{Solução}
  \begin{align*}
    \onslide<+->{A & = \int_0^b f(x)\,dx    \\[2mm]}
    \onslide<+->{& = \int_0^b x + 2 \,dx    \\[2mm]}
    \onslide<+->{& = \int_0^b x\, dx + 2 \int_0^b dx
                 & & \text{(Linearidade)}   \\[2mm]}
    \onslide<+->{& = \frac{b^2}{2} + 2b
                 & & \text{(Área do triângulo e do retângulo)}}
  \end{align*}
\end{resolution}

%------------------------------------------------------------------------------%

%------------------------------------------------------------------------------%
\section{Soma de Riemman}

%------------------------------------------------------------------------------%
\begin{frame}{Área entre o gráfico de $f$ e o eixo $x$ no intervalo $[a, b]$}
  \centering
  \bigskip
  \input{fig_integral_definida_area}
\end{frame}

%------------------------------------------------------------------------------%
\begin{frame}{Partição do intervalo $[a, b]$}
  \centering
  \bigskip
  % Title: Figure 1
% Creator: GL2PS 1.4.2, (C) 1999-2020 C. Geuzaine
% For: Octave
% CreationDate: Fri Aug  4 13:37:02 2023
\setlength{\unitlength}{1pt}
\begin{picture}(0,0)
\includegraphics[width=14cm,height=5cm]{fig_soma_riemman_particao}
\end{picture}%
\begin{picture}(396,141)(0,0)
\fontsize{10}{0}\selectfont\put(311.833,109.399){\makebox(0,0)[l]{\textcolor[rgb]{0,0,0}{{$y=f(x)$}}}}
\fontsize{10}{0}\selectfont\put(93.5842,8.22396){\makebox(0,0)[b]{\textcolor[rgb]{0,0,0}{{$a=x_0$}}}}
\fontsize{10}{0}\selectfont\put(121.956,8.22396){\makebox(0,0)[b]{\textcolor[rgb]{0,0,0}{{$x_1$}}}}
\fontsize{10}{0}\selectfont\put(147.963,8.22396){\makebox(0,0)[b]{\textcolor[rgb]{0,0,0}{{$x_2$}}}}
\fontsize{10}{0}\selectfont\put(173.969,8.22396){\makebox(0,0)[b]{\textcolor[rgb]{0,0,0}{{$x_3$}}}}
\fontsize{10}{0}\selectfont\put(199.976,8.22396){\makebox(0,0)[b]{\textcolor[rgb]{0,0,0}{{$x_4$}}}}
\fontsize{10}{0}\selectfont\put(224.942,8.22396){\makebox(0,0)[b]{\textcolor[rgb]{0,0,0}{{$x_5$}}}}
\fontsize{10}{0}\selectfont\put(250.948,8.22396){\makebox(0,0)[b]{\textcolor[rgb]{0,0,0}{{$\cdots$}}}}
\fontsize{10}{0}\selectfont\put(276.955,8.22396){\makebox(0,0)[b]{\textcolor[rgb]{0,0,0}{{$x_{n-1}$}}}}
\fontsize{10}{0}\selectfont\put(307.693,8.22396){\makebox(0,0)[b]{\textcolor[rgb]{0,0,0}{{$b=x_n$}}}}
\end{picture}

\end{frame}

%------------------------------------------------------------------------------%
\begin{frame}{Aproximação por retângulos}
  \centering
  \bigskip
  % Title: Figure 1
% Creator: GL2PS 1.4.2, (C) 1999-2020 C. Geuzaine
% For: Octave
% CreationDate: Fri Aug  4 13:37:03 2023
\setlength{\unitlength}{1pt}
\begin{picture}(0,0)
\includegraphics[width=14cm,height=5cm]{fig_soma_riemman_retangulos}
\end{picture}%
\begin{picture}(396,141)(0,0)
\fontsize{10}{0}\selectfont\put(311.833,109.399){\makebox(0,0)[l]{\textcolor[rgb]{0,0,0}{{$y=f(x)$}}}}
\fontsize{10}{0}\selectfont\put(93.5842,8.22396){\makebox(0,0)[b]{\textcolor[rgb]{0,0,0}{{$a=x_0$}}}}
\fontsize{10}{0}\selectfont\put(121.956,8.22396){\makebox(0,0)[b]{\textcolor[rgb]{0,0,0}{{$x_1$}}}}
\fontsize{10}{0}\selectfont\put(147.963,8.22396){\makebox(0,0)[b]{\textcolor[rgb]{0,0,0}{{$x_2$}}}}
\fontsize{10}{0}\selectfont\put(173.969,8.22396){\makebox(0,0)[b]{\textcolor[rgb]{0,0,0}{{$x_3$}}}}
\fontsize{10}{0}\selectfont\put(199.976,8.22396){\makebox(0,0)[b]{\textcolor[rgb]{0,0,0}{{$x_4$}}}}
\fontsize{10}{0}\selectfont\put(224.942,8.22396){\makebox(0,0)[b]{\textcolor[rgb]{0,0,0}{{$x_5$}}}}
\fontsize{10}{0}\selectfont\put(250.948,8.22396){\makebox(0,0)[b]{\textcolor[rgb]{0,0,0}{{$\cdots$}}}}
\fontsize{10}{0}\selectfont\put(276.955,8.22396){\makebox(0,0)[b]{\textcolor[rgb]{0,0,0}{{$x_{n-1}$}}}}
\fontsize{10}{0}\selectfont\put(307.693,8.22396){\makebox(0,0)[b]{\textcolor[rgb]{0,0,0}{{$b=x_n$}}}}
\fontsize{10}{0}\selectfont\put(99.4988,35.3245){\makebox(0,0)[bl]{\textcolor[rgb]{0,0,0}{{$A_1$}}}}
\fontsize{10}{0}\selectfont\put(127.871,35.3245){\makebox(0,0)[bl]{\textcolor[rgb]{0,0,0}{{$A_2$}}}}
\fontsize{10}{0}\selectfont\put(153.877,35.3245){\makebox(0,0)[bl]{\textcolor[rgb]{0,0,0}{{$A_3$}}}}
\fontsize{10}{0}\selectfont\put(179.884,35.3245){\makebox(0,0)[bl]{\textcolor[rgb]{0,0,0}{{$A_4$}}}}
\fontsize{10}{0}\selectfont\put(205.89,35.3245){\makebox(0,0)[bl]{\textcolor[rgb]{0,0,0}{{$A_5$}}}}
\fontsize{10}{0}\selectfont\put(282.869,35.3245){\makebox(0,0)[bl]{\textcolor[rgb]{0,0,0}{{$A_n$}}}}
\end{picture}

\end{frame}

%------------------------------------------------------------------------------%
\begin{frame}{Aproximação por retângulos}
\begin{align*}
  \onslide<+->{\Delta x & = \frac{b-a}{n}      \\[2mm]}
  \onslide<+->{x_k & = a + k\Delta x           \\[3mm]}
  \onslide<+->{A_k & = f(x_k)\Delta x          \\[6mm]}
  \onslide<+->{S_n & =   \sum_{k=1}^n  A_k }
  \onslide<+->{    \;=\; \sum_{k=1}^n f(x_k)\Delta x}
\end{align*}
\end{frame}

%------------------------------------------------------------------------------%
\begin{frame}{Área via soma de Riemann}
  A \emph{área} da região sob o gráfico de uma \emph{função contínua} $f$
  no intervalo $[a, b]$\\
  é dada pelo limite
  \[
    A = \lim_{n \to \infty} \; \sum_{i=1}^n f(x_i) \Delta x
  \]
\end{frame}

%------------------------------------------------------------------------------%
\begin{frame}{Definição Integral Definida -- Soma de Riemann}
  A \emph{Integral Definida} da função \(f\) no intervalo \([a, b]\) é

  \[
    \int_a^b f(x) dx
    \;=\; \lim_{n \to \infty} \; \sum_{i=1}^n f({\color<2->{magenta}c_i}) \Delta x
  \]

  \pause
  ${\color{magenta}c_i}\in[x_{i-1}, x_i]$ são os \emph{pontos amostrais}

  \pause
  \bigskip
  Desde que o limite exista

\end{frame}

%------------------------------------------------------------------------------%
\begin{frame}{Consequência da Definição}
  A área tem o mesmo \emph{sinal} da função

  \begin{center}
    \includegraphics[height=5cm]{integral_definida_03}
  \end{center}
\end{frame}

%------------------------------------------------------------------------------%
\section{Propriedades da integral definida}

%------------------------------------------------------------------------------%
\begin{frame}{Sentido de integração}
  \large
  \(
    \int_{\textcolor{blue}{a}}^{\textcolor{magenta}{b}} f(x)\,dx
    = \textcolor{red}{-}
    \int^{\textcolor{blue}{a}}_{\textcolor{magenta}{b}} f(x)\,dx
  \)
\end{frame}

%------------------------------------------------------------------------------%
\begin{frame}{Intervalo de comprimento nulo}
  \large
  \(
    \int_{\textcolor{blue}{a}}^{\textcolor{blue}{a}\vphantom{b}} f(x)dx
    = \textcolor{blue}{0}
  \)
\end{frame}

%------------------------------------------------------------------------------%
\begin{frame}{Área de um retângulo}{}
  \large
  \(
    \int_{a}^b {\color{blue}c}\,dx = {\color{blue}c}(b-a)
  \)
  \qquad\qquad
  para {\color{blue}$c$} constante
\end{frame}

%------------------------------------------------------------------------------%
\begin{frame}{Soma e subtração de funções}
  \large
  \(
    \int_{a}^b f(x) \;\memph{+}\; g(x)\, dx =
    \int_{a}^b f(x)\,dx \;\memph{+}\,
    \int_{a}^b g(x)\,dx
  \)
  \\[7mm]
  \(
    \int_{a}^b f(x) \;\memph{-}\; g(x)\, dx =
    \int_{a}^b f(x)\,dx \;\memph{-}\,
    \int_{a}^b g(x)\,dx
  \)
\end{frame}

%------------------------------------------------------------------------------%
\begin{frame}{Multiplicação de função por constante}{}
  \large
  \(
        \int_{a}^b \textcolor{blue}{c}\,f(x)\,dx
      = \textcolor{blue}{c} \int_{a}^b f(x)\,dx
  \)
  \qquad\qquad
  para \textcolor{blue}{$c$} constante
\end{frame}

%------------------------------------------------------------------------------%
\begin{frame}{Separando o intervalo em dois}{}
  \large
  \(
    \int_{\textcolor{blue}{a}}^{\textcolor{blue}{b}} f(x)\,dx
    = \int_{\textcolor{blue}{a}}^{\textcolor{magenta}{c}} f(x)\,dx
    + \int_{\textcolor{magenta}{c}}^{\textcolor{blue}{b}} f(x)\,dx
  \)
\end{frame}

%------------------------------------------------------------------------------%
\begin{frame}{Área positiva}
  \large
  Se \(f(x)\textcolor{blue}{\,\geq\,} 0\) para \(a \leq x \leq b\), então
  \(
    \int_{a}^b f(x)\,dx \textcolor{blue}{\,\geq\,} 0
  \)
\end{frame}

%------------------------------------------------------------------------------%
\begin{frame}{Comparação entre áreas}
  \large
  Se \textcolor{blue}{$f(x) \geq g(x)$}
  para $a \leq x \leq b$, então \\[9mm]
  \[
    \int_{a}^b f(x)\,dx \;\textcolor{blue}{\,\geq\,}\; \int_{a}^b g(x)\,dx
  \]
\end{frame}

%------------------------------------------------------------------------------%
\begin{frame}{Limitação inferior e superior}
  \large
  Se \({\color{blue}m} \leq f(x) \leq {\color{red}M}\) para $a \leq x \leq b$ então \\[9mm]
  \[
    {\color{blue} m}(b-a)  \quad\leq\quad
    \int_{a}^b f(x)\,dx    \quad\leq\quad
    {\color{red} M}(b-a)
  \]
\end{frame}

%------------------------------------------------------------------------------%
\section{Exemplos}

%------------------------------------------------------------------------------%
\begin{question}
  Calcule \(\int_{-7}^7 x^3 dx\)
\end{question}

%------------------------------------------------------------------------------%
\begin{resolution}{Solução}
  \[
    A
    = \int_{-7}^7 x^3 dx
    \pause
    = \int_{-7}^0 x^3 dx + \int_0^7 x^3 dx
  \]
  \pause
  por simetria
  \[
    \int_{-7}^0 x^3 dx = - \int_0^7 x^3 dx
  \]
  \pause
  portanto
  \[
    A = 0
  \]
\end{resolution}

%------------------------------------------------------------------------------%
%------------------------------------------------------------------------------%
\begin{question}
  Calcule \(\int_{-2}^4 \abs{x} \,dx\)
\end{question}

%------------------------------------------------------------------------------%
\begin{resolution}{Solução}
  \begin{align*}
    \onslide<+->{A
    & = \int_{-2}^4 \abs{x} \,dx                                       \\[1mm]}
    \onslide<+->{& = \int_{-2}^0 \abs{x} \,dx + \int_0^4 \abs{x} \,dx  \\[1mm]}
    \onslide<+->{& = \int_{-2}^0 -x \,dx + \int_0^4 x \,dx             \\[1mm]}
    \onslide<+->{& = \frac{2^2}{2} + \frac{4^2}{2}                     \\[1mm]}
    \onslide<+->{& = 2 + 8                                             \\[1mm]}
    \onslide<+->{& = 10}
  \end{align*}
\end{resolution}

%------------------------------------------------------------------------------%

%------------------------------------------------------------------------------%
\begin{question}
  Calcule a integral da função de Heaviside
  \[
  \heaviside(x) = \begin{cases}
    \; 0 & \; x < 0 \\
    \; 1 & \; x \geq 0
  \end{cases}
  \]
  no intervalo $[-1, 2]$.
\end{question}

%------------------------------------------------------------------------------%
\begin{resolution}{Solução}
  \[
    A
    = \int_{-1}^{2} \heaviside(x) \,dx
    \pause
    = \int_{-1}^{0} 0 \,dx
    + \int_{0}^{2}  1 \,dx
  \]
  \pause
  \[
    \int_{-1}^{0} 0 \,dx
    \pause
    = 0
  \]
  \pause
  \[
    \int_{0}^{2} 1 \,dx
    \pause
    \;=\; \text{base} \times \text{altura}
    \pause
    \;=\; (2-0)\times 1
    \pause
    \;=\; 2
  \]
  \pause
  \[
    A
    \pause
    = 0 + 2
    \pause
    = 2
  \]
\end{resolution}

%------------------------------------------------------------------------------%
\input{ex/B-02-Integrais_Definidas_ex-5}

%------------------------------------------------------------------------------%
\section{Lista Mínima}

%------------------------------------------------------------------------------%
\begin{frame}{Lista Mínima}
  Estudar a Seção 2.3 da Apostila

  Exercícios: 4, 5

  \vfill
  Atenção: A prova é baseada no livro, não nas apresentações
\end{frame}

%------------------------------------------------------------------------------%
\end{document}
%------------------------------------------------------------------------------%
